\documentclass[12pt,a4paper]{article}
\usepackage[UTF8]{ctex}
\usepackage{amsmath,amssymb,amsthm}
\usepackage{geometry}

\geometry{margin=2.5cm}

\title{相对角速度$\omega_{\text{rel}}$的含义和表示}
\author{第8章补充说明}
\date{}

\begin{document}

\maketitle

\section{问题提出}

\textbf{问题}：如何理解"$\omega_{\text{rel}} \in \mathbb{R}^3$：坐标系$\{i\}$相对于$\{i-1\}$的相对角速度（在$\{i\}$中表示）"？

\section{基本概念}

\subsection{相对角速度的定义}

\textbf{定义}：$\omega_{\text{rel}}$是坐标系$\{i\}$相对于坐标系$\{i-1\}$的角速度，也可以理解为坐标系$\{i\}$自身的角速度。

\textbf{物理意义}：
\begin{itemize}
\item 如果坐标系$\{i-1\}$固定不动，$\omega_{\text{rel}}$描述了坐标系$\{i\}$自身的旋转运动
\item 这是坐标系$\{i\}$自身的角速度，由关节$i$的运动产生
\item 这是两个坐标系之间的相对旋转运动
\end{itemize}

\textbf{关键理解}：
\begin{itemize}
\item $\omega_{\text{rel}}$是坐标系$\{i\}$\textbf{自身}的角速度
\item 它描述了坐标系$\{i\}$相对于$\{i-1\}$的旋转
\item 由关节$i$的运动直接产生：$\omega_{\text{rel}} = \dot{\theta}_i z_i$
\end{itemize}

\subsection{表示坐标系的选择}

\textbf{关键问题}：$\omega_{\text{rel}}$应该在哪个坐标系中表示？

\textbf{答案}：在坐标系$\{i\}$中表示。

\textbf{原因}：
\begin{enumerate}
\item 关节轴$z_i$通常在坐标系$\{i\}$中定义
\item 因此，相对角速度$\omega_{\text{rel}} = \dot{\theta}_i z_i$自然在$\{i\}$中表示
\item 这样便于在递归算法中使用
\end{enumerate}

\section{数学表达}

\subsection{一般形式}

\textbf{相对角速度的表达式}：
\begin{equation}
\omega_{\text{rel}} = \dot{\theta}_i z_i
\end{equation}

\textbf{其中}：
\begin{itemize}
\item $\dot{\theta}_i$：关节$i$的角速度（标量）
\item $z_i$：关节$i$的旋转轴（在坐标系$\{i\}$中表示的单位向量）
\item $\omega_{\text{rel}}$：相对角速度（在坐标系$\{i\}$中表示的3维向量）
\end{itemize}

\subsection{坐标表示}

\textbf{在坐标系$\{i\}$中的表示}：

如果$z_i$在坐标系$\{i\}$中的表示为：
\begin{equation}
z_i = \begin{bmatrix} z_{i,x} \\ z_{i,y} \\ z_{i,z} \end{bmatrix}_{\{i\}}
\end{equation}

那么$\omega_{\text{rel}}$在坐标系$\{i\}$中的表示为：
\begin{equation}
\omega_{\text{rel}} = \dot{\theta}_i \begin{bmatrix} z_{i,x} \\ z_{i,y} \\ z_{i,z} \end{bmatrix}_{\{i\}} = \begin{bmatrix} \dot{\theta}_i z_{i,x} \\ \dot{\theta}_i z_{i,y} \\ \dot{\theta}_i z_{i,z} \end{bmatrix}_{\{i\}}
\end{equation}

\textbf{注意}：下标$\{i\}$表示这个向量在坐标系$\{i\}$中的坐标表示。

\section{具体例子}

\subsection{例子1：绕$z$轴旋转}

\textbf{场景}：

考虑一个旋转关节，关节轴$z_i$与坐标系$\{i\}$的$z$轴对齐：
\begin{equation}
z_i = \begin{bmatrix} 0 \\ 0 \\ 1 \end{bmatrix}_{\{i\}}
\end{equation}

关节以角速度$\dot{\theta}_i = 2$ rad/s旋转。

\textbf{自身角速度}：

\begin{equation}
\omega_{\text{rel}} = \dot{\theta}_i z_i = 2 \times \begin{bmatrix} 0 \\ 0 \\ 1 \end{bmatrix}_{\{i\}} = \begin{bmatrix} 0 \\ 0 \\ 2 \end{bmatrix}_{\{i\}}
\end{equation}

\textbf{物理意义}：
\begin{itemize}
\item 坐标系$\{i\}$自身的角速度为$[0, 0, 2]^T$ rad/s（在$\{i\}$中表示）
\item 坐标系$\{i\}$绕自身的$z$轴旋转
\item 角速度大小为$2$ rad/s
\item 方向沿$z$轴正方向（在$\{i\}$中）
\end{itemize}

\textbf{验证}：

在坐标系$\{i\}$中，这个向量表示：
\begin{itemize}
\item $x$方向分量：$0$ rad/s
\item $y$方向分量：$0$ rad/s
\item $z$方向分量：$2$ rad/s
\end{itemize}

这符合绕$z$轴旋转的物理意义。

\subsection{例子2：绕任意轴旋转}

\textbf{场景}：

考虑一个旋转关节，关节轴$z_i$在坐标系$\{i\}$中的表示为：
\begin{equation}
z_i = \frac{1}{\sqrt{3}} \begin{bmatrix} 1 \\ 1 \\ 1 \end{bmatrix}_{\{i\}}
\end{equation}

这是一个单位向量，指向$(1,1,1)$方向。

关节以角速度$\dot{\theta}_i = \pi$ rad/s旋转。

\textbf{自身角速度}：

\begin{equation}
\omega_{\text{rel}} = \dot{\theta}_i z_i = \pi \times \frac{1}{\sqrt{3}} \begin{bmatrix} 1 \\ 1 \\ 1 \end{bmatrix}_{\{i\}} = \frac{\pi}{\sqrt{3}} \begin{bmatrix} 1 \\ 1 \\ 1 \end{bmatrix}_{\{i\}}
\end{equation}

\textbf{物理意义}：
\begin{itemize}
\item 坐标系$\{i\}$自身的角速度为$\frac{\pi}{\sqrt{3}}[1, 1, 1]^T$ rad/s（在$\{i\}$中表示）
\item 坐标系$\{i\}$绕自身的轴$(1,1,1)$旋转
\item 角速度大小为$\pi$ rad/s
\item 方向沿$(1,1,1)$方向（在$\{i\}$中）
\end{itemize}

\textbf{验证}：

在坐标系$\{i\}$中，这个向量表示：
\begin{itemize}
\item $x$方向分量：$\frac{\pi}{\sqrt{3}}$ rad/s
\item $y$方向分量：$\frac{\pi}{\sqrt{3}}$ rad/s
\item $z$方向分量：$\frac{\pi}{\sqrt{3}}$ rad/s
\end{itemize}

所有分量相等，符合绕$(1,1,1)$方向旋转的物理意义。

\subsection{例子3：在机器人手臂中的应用}

\textbf{场景}：

考虑一个6自由度机器人手臂的第3个关节：
\begin{itemize}
\item 关节类型：旋转关节
\item 关节轴$z_3$在坐标系$\{3\}$中的表示：$z_3 = [0, 0, 1]^T$（沿$z$轴）
\item 关节角速度：$\dot{\theta}_3 = 1.5$ rad/s
\end{itemize}

\textbf{自身角速度}：

\begin{equation}
\omega_{\text{rel}} = \dot{\theta}_3 z_3 = 1.5 \times \begin{bmatrix} 0 \\ 0 \\ 1 \end{bmatrix}_{\{3\}} = \begin{bmatrix} 0 \\ 0 \\ 1.5 \end{bmatrix}_{\{3\}}
\end{equation}

\textbf{物理意义}：
\begin{itemize}
\item 坐标系$\{3\}$自身的角速度为$[0, 0, 1.5]^T$ rad/s（在$\{3\}$中表示）
\item 坐标系$\{3\}$绕自身的$z$轴旋转
\item 角速度大小为$1.5$ rad/s
\item 这个角速度由关节3的运动产生
\end{itemize}

\textbf{在速度合成中的应用}：

如果坐标系$\{2\}$的角速度为$\omega_2 = [0.5, 0, 0]^T$（在$\{2\}$中表示），那么坐标系$\{3\}$的绝对角速度为：

\begin{align}
\omega_3 &= R_3^T \omega_2 + \omega_{\text{rel}} \\
&= R_3^T \begin{bmatrix} 0.5 \\ 0 \\ 0 \end{bmatrix}_{\{2\}} + \begin{bmatrix} 0 \\ 0 \\ 1.5 \end{bmatrix}_{\{3\}}
\end{align}

\textbf{注意}：
\begin{itemize}
\item $R_3^T \omega_2$将$\omega_2$从$\{2\}$转换到$\{3\}$
\item $\omega_{\text{rel}}$已经在$\{3\}$中，可以直接相加
\item 所有项都在$\{3\}$中表示，可以相加
\end{itemize}

\section{为什么在$\{i\}$中表示？}

\subsection{原因1：与关节轴定义一致}

\textbf{关节轴的定义}：

在机器人学中，关节轴$z_i$通常在坐标系$\{i\}$中定义：
\begin{equation}
z_i = \begin{bmatrix} z_{i,x} \\ z_{i,y} \\ z_{i,z} \end{bmatrix}_{\{i\}}
\end{equation}

\textbf{自身角速度的定义}：

坐标系$\{i\}$自身的角速度是沿关节轴方向的：
\begin{equation}
\omega_{\text{rel}} = \dot{\theta}_i z_i
\end{equation}

\textbf{结论}：

由于$z_i$在$\{i\}$中定义，$\omega_{\text{rel}} = \dot{\theta}_i z_i$自然也在$\{i\}$中表示。这反映了坐标系$\{i\}$自身的旋转运动。

\subsection{原因2：便于递归计算}

\textbf{递归算法的特点}：

在递归牛顿-欧拉算法中，所有量都在当前坐标系$\{i\}$中表示：
\begin{itemize}
\item $\omega_i$：在$\{i\}$中表示
\item $v_i$：在$\{i\}$中表示
\item $\omega_{\text{rel}}$：在$\{i\}$中表示
\item $v_{\text{rel}}$：在$\{i\}$中表示
\end{itemize}

\textbf{速度合成公式}：

\begin{align}
\omega_i &= R_i^T \omega_{i-1} + \omega_{\text{rel}} \\
v_i &= R_i^T (v_{i-1} + \omega_{i-1} \times p_i) + v_{\text{rel}}
\end{align}

\textbf{关键点}：

\begin{itemize}
\item $R_i^T \omega_{i-1}$将$\omega_{i-1}$从$\{i-1\}$转换到$\{i\}$
\item $\omega_{\text{rel}}$已经在$\{i\}$中
\item 两者都在$\{i\}$中，可以直接相加
\end{itemize}

\subsection{原因3：物理直观性}

\textbf{物理意义}：

$\omega_{\text{rel}}$是坐标系$\{i\}$自身的角速度，描述了坐标系$\{i\}$自身的旋转运动。

\textbf{在$\{i\}$中表示的优势}：

\begin{itemize}
\item 直接反映了坐标系$\{i\}$自身的旋转运动
\item 与坐标系$\{i\}$的基向量对齐，物理意义清晰
\item 便于理解坐标系$\{i\}$自身的运动状态
\item 如果坐标系$\{i-1\}$固定，$\omega_{\text{rel}}$就是坐标系$\{i\}$的绝对角速度
\end{itemize}

\section{与其他表示方法的对比}

\subsection{方法1：在$\{i\}$中表示（标准方法）}

\textbf{表达式}：
\begin{equation}
\omega_{\text{rel}} = \dot{\theta}_i z_i \quad \text{（在$\{i\}$中表示）}
\end{equation}

\textbf{优点}：
\begin{itemize}
\item 与关节轴定义一致
\item 便于递归计算
\item 物理意义直观
\end{itemize}

\textbf{速度合成}：
\begin{equation}
\omega_i = R_i^T \omega_{i-1} + \omega_{\text{rel}}
\end{equation}

\subsection{方法2：在$\{i-1\}$中表示（理论可能）}

\textbf{表达式}：
\begin{equation}
\omega_{\text{rel},\{i-1\}} = R_i \omega_{\text{rel},\{i\}} = R_i (\dot{\theta}_i z_i)
\end{equation}

\textbf{缺点}：
\begin{itemize}
\item 需要额外的坐标变换
\item 与关节轴定义不一致
\item 在递归算法中不常用
\end{itemize}

\textbf{速度合成}（如果使用这种方法）：
\begin{equation}
\omega_i = R_i^T (\omega_{i-1} + \omega_{\text{rel},\{i-1\}})
\end{equation}

\textbf{注意}：这种方法在理论上可行，但在实际应用中不常用。

\section{总结}

\subsection{关键要点}

\begin{enumerate}
\item \textbf{定义}：
\begin{itemize}
\item $\omega_{\text{rel}}$是坐标系$\{i\}$自身的角速度
\item 它描述了坐标系$\{i\}$相对于$\{i-1\}$的旋转
\item 表达式：$\omega_{\text{rel}} = \dot{\theta}_i z_i$
\end{itemize}

\item \textbf{表示}：
\begin{itemize}
\item $\omega_{\text{rel}}$在坐标系$\{i\}$中表示
\item 原因：与关节轴$z_i$的定义一致，反映坐标系$\{i\}$自身的旋转
\end{itemize}

\item \textbf{物理意义}：
\begin{itemize}
\item 这是坐标系$\{i\}$自身的角速度，由关节$i$的运动产生
\item 如果坐标系$\{i-1\}$固定，$\omega_{\text{rel}}$就是坐标系$\{i\}$的绝对角速度
\item 方向沿关节轴$z_i$，大小等于$\dot{\theta}_i$
\end{itemize}

\item \textbf{应用}：
\begin{itemize}
\item 在速度合成公式中使用：$\omega_i = R_i^T \omega_{i-1} + \omega_{\text{rel}}$
\item 所有项都在$\{i\}$中表示，可以直接相加
\end{itemize}
\end{enumerate}

\subsection{公式总结}

\textbf{自身角速度}：
\begin{equation}
\omega_{\text{rel}} = \dot{\theta}_i z_i \quad \text{（在$\{i\}$中表示）}
\end{equation}

\textbf{其中}：
\begin{itemize}
\item $\dot{\theta}_i$：关节$i$的角速度（标量）
\item $z_i$：关节$i$的旋转轴（在$\{i\}$中表示的单位向量）
\item $\omega_{\text{rel}} \in \mathbb{R}^3$：坐标系$\{i\}$自身的角速度（在$\{i\}$中表示的3维向量）
\end{itemize}

\textbf{关键理解}：
\begin{itemize}
\item $\omega_{\text{rel}}$是坐标系$\{i\}$自身的角速度
\item 它由关节$i$的运动直接产生
\item 在坐标系$\{i\}$中表示，反映了坐标系$\{i\}$自身的旋转状态
\end{itemize}

\end{document}

